\\section{Model for Problem 2}

\\subsection{Data and continuous reconstruction}
For heat $h$, the source time is defined by $t_{h,i}=0.5(s_{h,i}-s_{h,1})$ seconds. Let $O_h(t)$ denote cumulative oxygen consumption in $\\mathrm{m^3}$ and let $C_{g,h}(t)$ denote the instantaneous volumetric flow rate of gas $g\\in\\{\\mathrm{CO},\\mathrm{CO_2}\\}$ in $\\mathrm{m^3/s}$. A robust Hampel detector with a zero-scale fallback identifies isolated spikes while preserving sustained steps. PCHIP is then fitted to the retained observations and the original values, cleaned values, and masks are saved separately.

\\subsection{Oxygen flow and cumulative emissions}
The oxygen flow is estimated from the reconstructed cumulative oxygen curve using a third-order, 31-sample Savitzky--Golay derivative,
\\[
\\widehat q_h(t)=\\max\\{0,\\operatorname{SGDeriv}_{31,3}[\\widehat O_h](t)\\}.
\\]
The cumulative gas emission is computed by
\\[
\\widehat E_{g,h}(t)=\\int_0^t\\widehat C_{g,h}(u)\\,du.
\\]
PCHIP analytic derivatives and antiderivatives are retained as consistency diagnostics. The analytic integral recovers the PCHIP curve to machine precision, whereas trapezoidal integration of nodal derivatives is unsuitable for the long plateaus and discrete changes in the cumulative record.

\\subsection{Oxygen--flue-gas relationships}
For each heat, four relationships are evaluated: cumulative oxygen versus instantaneous CO, instantaneous CO$_2$, cumulative CO emission, and cumulative CO$_2$ emission. Pearson and Spearman coefficients, linear and quadratic fits, and oxygen-quantile stage statistics are reported. High correlations between cumulative quantities are interpreted cautiously because both variables share a time trend.

\\subsection{Reliability-weighted local delay}
Let $x_h(t)=\\widehat q_h(t)$ and $y_{g,h}(t)=\\widehat C_{g,h}(t)$. In each window, detrended changes are compared over nonnegative lags. A window is accepted only when source and target dynamic standard deviations, the maximum correlation, peak separation, peak width, and boundary checks all pass. Thus
\\[
\\mathcal V_g=\\{i:\\text{window }i\\text{ passes all reliability checks}\\},
\\]
and $\\widehat\\tau_i$ is stored only for $i\\in\\mathcal V_g$. Unidentifiable windows remain NaN and are accompanied by validity masks and confidence scores.

\\subsection{Flow-dependent delay and time mapping}
Reliable window samples are $(\\bar q_i,\\widehat\\tau_i,w_i)$. We compare a fixed delay with the constrained inverse-square-root model
\\[
\\tau(q)=a+bq^{-1/2},\\qquad a\\geq0,\\quad b\\geq0.
\\]
The model is fitted only when at least five reliable samples are available. The source-to-measurement map is
\\[
\\phi_g(t)=t+\\tau_g(t),
\\]
and only in-domain, strictly increasing observed mappings are used for offline alignment. Each heat file contains observed delay, model delay, confidence, observed and model mappings, and alignment masks.

\\subsection{Results and limitations}
All ten heats have continuous Sequence and a 0.5-s sampling interval. Reliable delay coverage is low (0--4.7\\%), and many leave-one-heat-out test heats have no valid delay windows. The resulting delay is a statistical response delay containing reaction, transport, and instrument components; it is not an independently measured pipe transport time. The NPZ interface preserves NaNs in unidentifiable intervals so that later models can use the masks explicitly.
